\documentclass[a4paper,12pt]{article}
\usepackage[T2A]{fontenc}
\usepackage[utf8]{inputenc}
\usepackage[russian]{babel}
\usepackage{amsmath,amssymb,mathtools}
\usepackage{geometry}
\geometry{left=23mm,right=23mm,top=24mm,bottom=24mm}
\usepackage{microtype}
\usepackage{icomma}
\usepackage{tikz}
\usetikzlibrary{arrows.meta,positioning,shapes.geometric,calc}
\usepackage{pgfplots}
\pgfplotsset{compat=1.18}
\usepackage{booktabs,tabularx,array}
\usepackage{enumitem}
\usepackage{fancyhdr}
\usepackage{hyperref}
\usepackage{xcolor}
\definecolor{accent}{HTML}{007C84}
\definecolor{ink}{HTML}{283039}
\definecolor{muted}{HTML}{617078}
\hypersetup{colorlinks=true,linkcolor=accent,urlcolor=accent}
\setlength{\parindent}{0pt}
\setlength{\parskip}{5pt}
\setlist[itemize]{leftmargin=1.7em,itemsep=3pt,topsep=4pt}
\setlist[enumerate]{leftmargin=1.7em,itemsep=10pt,topsep=6pt}
\pagestyle{fancy}
\fancyhf{}
\fancyhead[L]{\small\color{muted}Графики и неравенства}
\fancyhead[R]{\small\color{muted}Екатерина Скелли}
\fancyfoot[C]{\small\color{muted}\thepage}
\renewcommand{\headrulewidth}{0.4pt}
\setlength{\headheight}{15pt}
\newcommand{\keyidea}[1]{\textcolor{accent}{\textbf{#1}}}
\newcommand{\lessonsection}[1]{\section*{#1}}

\begin{document}
\begin{titlepage}
\thispagestyle{empty}
\pagenumbering{gobble}
\centering
\vspace*{15mm}
{\color{accent}\large МАТЕМАТИКА\par}
\vspace{15mm}
{\fontsize{29}{34}\selectfont Чек-лист\par}
\vspace{15mm}
{\color{accent}\LARGE Графики функций\\и линейные неравенства\par}
\vspace{12mm}
{\large Чтение графиков, преобразования\\и текстовые задачи\par}
\vspace{20mm}
{\large Екатерина Скелли\par}
\vspace{10mm}
{28 октября 2026 г.\par}
\vfill
{\small Лёвин Артём Александрович\\эксклюзивно для Екатерины Скелли\par}
\vspace{10mm}
\begin{tikzpicture}[x=1cm,y=1cm]
  \draw[accent,line width=.8pt]
     (0,0) .. controls (0.7,.4) and (1.6,-.4) .. (2.5,0)
           .. controls (3.4,.4) and (4.3,-.4) .. (5.1,0);
\end{tikzpicture}
\end{titlepage}
\pagenumbering{arabic}
\setcounter{page}{2}

\lessonsection{1. Как определить знак функции по графику}
\keyidea{Ключевая мысль.} Значение функции при данном $x$ --- это координата $y$ соответствующей точки графика.
\begin{itemize}
\item График \textbf{выше} оси $Ox$: $f(x)>0$.
\item График \textbf{ниже} оси $Ox$: $f(x)<0$.
\item График \textbf{пересекает или касается} оси $Ox$: $f(x)=0$ в соответствующих точках.
\end{itemize}
\keyidea{Как читать условие «при всех $x<1$ функция положительна»?}
Рассматривайте только часть графика \emph{левее вертикали $x=1$}. Чтобы утверждение было верным, вся эта часть должна находиться выше оси $Ox$. Если хотя бы одна точка лежит на оси или ниже неё, утверждение неверно.

\textbf{Проверяйте:} интервалы записывают по оси $x$, а знак функции определяют по высоте точек, то есть по оси $y$.

\lessonsection{2. Как сравнивать значение функции с числом}
Если нужно определить, где $f(x)<c$, мысленно проведите горизонтальную прямую $y=c$ и найдите участки графика \textbf{ниже} неё. Для $f(x)>c$ ищите участки \textbf{выше} этой прямой.

\begin{center}
\begin{tikzpicture}
\begin{axis}[
 width=.83\linewidth,height=5.1cm,
 axis lines=middle, xmin=-.3,xmax=5.15,ymin=-8.4,ymax=6.6,
 xtick={1,2,3,4,5},ytick={-8,-4,1,4},
 grid=major,grid style={gray!16},
 xlabel={$x$},ylabel={$y$},
 xlabel style={at={(axis description cs:1.02,.43)},anchor=west},
 ylabel style={at={(axis description cs:.055,1.01)},anchor=south},
 tick label style={font=\scriptsize},
 clip=false
]
\addplot[accent,thick,domain=0:5,samples=2]{3*x-8};
\addplot[muted,dashed,thin,domain=0:5]{1};
\addplot[only marks,mark=*,mark size=1.7pt,accent] coordinates {(3,1)};
\node[above right,font=\scriptsize] at (axis cs:3,1) {$(3;1)$};
\node[accent,above,font=\scriptsize] at (axis cs:4.35,5.4) {$y=3x-8$};
\end{axis}
\end{tikzpicture}
\end{center}

\textbf{Пример.} По графику функции $f(x)=3x-8$ определите, когда $f(x)<1$. Прямая $y=1$ пересекает график при $x=3$. Левее точки пересечения график ниже горизонтали. Поэтому
\[\boxed{x<3}.\]
При $x=3$ значение равно $1$, и строгий знак «меньше» исключает границу.

\clearpage
\lessonsection{3. Как перевести текст в знак неравенства}
\begin{center}
\renewcommand{\arraystretch}{1.35}
\begin{tabularx}{\linewidth}{X>{\centering\arraybackslash}p{5cm}}
\toprule
\textbf{Слова в условии}&\textbf{Математическая запись}\\
\midrule
«меньше», «ниже» & $<$\\
«не больше», «не превосходит» & $\le$\\
«больше», «выше» & $>$\\
«не меньше», «не менее» & $\ge$\\
\bottomrule
\end{tabularx}
\end{center}
Например, «$f(x)=-x+14$ не превосходит $2$» означает
\vspace{-7pt}
\[-x+14\le2\quad\Longleftrightarrow\quad -x\le-12\quad\Longleftrightarrow\quad x\ge12.\]
При делении на отрицательное число знак неравенства меняется на противоположный.

\lessonsection{4. Скобки защищают выражение}
\keyidea{Ключевая мысль.} Когда вычитают, умножают или подставляют целое выражение, удобно заключать его в скобки. Это предотвращает потерю знаков.

\textbf{Вычитание}
\[(2x+3)-(4x-7)=2x+3-4x+7=-2x+10.\]
\textbf{Умножение}
\[3(2-x)=6-3x,\qquad 3\cdot2-x=6-x.\]
\textbf{Подстановка.} Если $x=t-7$, то
\[2x+3=2(t-7)+3=2t-11.\]
\textbf{Дробное неравенство}
\[\frac{2(1-x)}{4}+\frac{12}{4}<\frac{12x}{4}-\frac{2x+1}{4}.\]
Умножаем обе части на положительное число $4$ и сохраняем скобки:
\[2(1-x)+12<12x-(2x+1).\]
\[14-2x<10x-1\quad\Longleftrightarrow\quad15<12x\quad\Longleftrightarrow\quad x>\frac54.\]
\textbf{Проверяйте:} после минуса перед скобкой меняются знаки всех слагаемых внутри неё.

\lessonsection{5. Числовая прямая и пересечение промежутков}
\begin{itemize}
\item Строгие знаки $<$ и $>$: граничная точка выколота.
\item Нестрогие знаки $\le$ и $\ge$: граничная точка включена.
\item «Решения из заданного промежутка» означает пересечение двух множеств.
\end{itemize}

\textbf{Пример.} Найдите решения неравенства $2(3x+5)-7<15$, принадлежащие $[-7;12]$.
\[6x+10-7<15\quad\Longleftrightarrow\quad6x<12\quad\Longleftrightarrow\quad x<2.\]
\[( -\infty;2)\cap[-7;12]=[-7;2).\]
\begin{center}
\begin{tikzpicture}[x=.42cm,y=1cm]
\draw[-{Stealth},gray!70] (-8,0)--(13,0);
\draw[accent,very thick] (-7,0)--(2,0);
\fill[accent] (-7,0) circle (2pt);
\draw[accent,fill=white,thick] (2,0) circle (2pt);
\foreach \x in {-7,-4,0,2,6,12} {\draw (\x,.08)--(\x,-.08) node[below,font=\scriptsize]{$\x$};}
\end{tikzpicture}
\end{center}
Если нужны только целые решения, перечисляем их:
\[-7,-6,-5,-4,-3,-2,-1,0,1.\]
Интервал и список целых чисел отвечают на разные вопросы.

\lessonsection{6. Текстовые задачи: составляем математическую модель}
\textbf{Перевод денег с комиссией.} Пусть $p$ --- сумма перевода в рублях. Комиссия: $12$ руб. и $4\%$ от перевода. Всего доступно $900$ руб. Тогда
\[p+12+0{,}04p\le900,\qquad 1{,}04p\le888,\qquad p\le\frac{888}{1{,}04}=\frac{11100}{13}.\]
Число $\frac{11100}{13}=853+\frac{11}{13}$ лежит между $853$ и $854$. Поэтому наибольшая целая сумма перевода равна \textbf{853 руб.}

\textbf{Проверяйте:} процент вычисляем от переводимой суммы $p$, а не от всего бюджета.

\textbf{Три последовательных чётных числа.} Пусть первое чётное натуральное число равно $x$. Следующие равны $x+2$ и $x+4$. Если сумма не больше $164$, то
\[x+(x+2)+(x+4)\le164,\qquad3x\le158,\qquad x\le\frac{158}{3}.\]
Наибольшее подходящее чётное $x=52$. Искомые числа: $52,54,56$ (их сумма $162$).

\clearpage
\lessonsection{7. Неравенства с любым числом решений и с параметром}
Иногда после упрощения переменная исчезает. Тогда проверяют оставшееся числовое утверждение.
\[
\begin{aligned}
3x-1>3x&\quad\Longleftrightarrow\quad -1>0&&\Rightarrow\ \text{решений нет};\\
x+4>x&\quad\Longleftrightarrow\quad 4>0&&\Rightarrow\ x\in\mathbb R.
\end{aligned}
\]
\textbf{Простейший параметр.} Для неравенства $2y+3a<2y+7$ получаем $3a<7$. Значит:
\begin{itemize}
\item если $a<\frac73$, подходит любое $y$;
\item если $a\ge\frac73$, решений по $y$ нет.
\end{itemize}
Параметр $a$ здесь задаёт фиксированное число. Например, при $a=-10$ неравенство верно для любого $y$, а при $a=10$ --- ни для одного.

\clearpage
\begin{center}
{\LARGE\color{accent} Алгоритм решения линейного неравенства\par}
\vspace{10mm}
\begin{tikzpicture}[
 >=Stealth,
 flow/.style={draw=accent,thin,align=center,font=\small,minimum height=9mm,inner sep=4pt},
 term/.style={flow,ellipse,minimum width=37mm},
 proc/.style={flow,rectangle,rounded corners=2pt,text width=67mm},
 decision/.style={flow,diamond,aspect=2.9,text width=17mm,inner sep=1pt},
 smallproc/.style={flow,rectangle,rounded corners=2pt,text width=36mm},
 arrow/.style={-{Stealth[length=2mm]},thin,draw=ink}
]
\node[term] (start) at (0,0) {Начало};
\node[proc] (simplify) at (0,-1.5) {Раскройте скобки, приведите подобные\\и запишите $Ax\mathrel{\square}B$};
\node[decision] (zero) at (0,-3.2) {$A=0$?};
\node[smallproc] (all) at (-4.5,-5.3) {Проверьте $0\mathrel{\square}B$.\\Если верно --- все $x$;\\если нет --- решений нет};
\node[decision] (pos) at (3,-4.65) {$A>0$?};
\node[smallproc] (keep) at (0.8,-6.65) {Разделите на $A$.\\Знак сохраняется};
\node[smallproc] (flip) at (5.0,-6.65) {Разделите на $A<0$.\\Знак меняется};
\node[proc] (write) at (0,-9.1) {Запишите множество решений;\\проверьте включение граничной точки};
\node[term] (end) at (0,-10.45) {Конец};
\draw[arrow] (start)--(simplify);
\draw[arrow] (simplify)--(zero);
\draw[arrow] (zero.west)--(-4.5,-3.2)--node[left,pos=.55,font=\scriptsize]{да}(all.north);
\draw[arrow] (zero.east)-|node[above,pos=.3,font=\scriptsize]{нет}(pos.north);
\draw[arrow] (pos.south west)-|node[above,pos=.3,font=\scriptsize]{да}(keep.north);
\draw[arrow] (pos.south east)-|node[above,pos=.3,font=\scriptsize]{нет}(flip.north);
\draw[arrow] (all.south)--(-4.5,-9.1)--(write.west);
\draw[arrow] (keep.south)--(0.8,-8.2)--(0,-8.2)--(write.north);
\draw[arrow] (flip.south)--(5.0,-9.1)--(write.east);
\draw[arrow] (write)--(end);
\end{tikzpicture}
\end{center}

\clearpage
\lessonsection{Тренировочный блок: Путь А}
\emph{10 заданий. Выполняйте последовательно: от чтения графика к текстовой математической модели. В заданиях с неравенствами указывайте множество решений.}
\begin{enumerate}[label=\textbf{\arabic*.}]
\item По графику определите, при каких $x$ выполняются $f(x)>0$, $f(x)=0$ и $f(x)<0$.
\begin{center}
\begin{tikzpicture}
\begin{axis}[width=.66\linewidth,height=4cm,axis lines=middle,xmin=-2.5,xmax=5.1,ymin=-3.5,ymax=5.1,
 xtick={-2,2,4},ytick={-2,2,4},grid=major,grid style={gray!15},xlabel={$x$},ylabel={$y$},tick label style={font=\scriptsize},clip=false]
\addplot[accent,thick,domain=-2.5:5.1]{2-x};
\end{axis}
\end{tikzpicture}
\end{center}
\item По графику прямой определите все $x$, при которых $f(x)\le1$. Используйте только графический способ.
\begin{center}
\begin{tikzpicture}
\begin{axis}[width=.66\linewidth,height=4cm,axis lines=middle,xmin=-1,xmax=5,ymin=-6.2,ymax=8.0,
 xtick={-1,1,2,3,4,5},ytick={-6,-4,-2,1,3,5,7},grid=major,grid style={gray!15},xlabel={$x$},ylabel={$y$},tick label style={font=\scriptsize},clip=false]
\addplot[accent,thick,domain=-1:5]{5-2*x};
\end{axis}
\end{tikzpicture}
\end{center}
\item Решите неравенство $2x+6\le0$. Изобразите ответ на числовой прямой.
\item При каких значениях $x$ функция $f(x)=-x+14$ \textbf{не превосходит} числа $2$? Запишите и решите соответствующее неравенство.
\end{enumerate}

\clearpage
\begin{enumerate}[label=\textbf{\arabic*.},start=5]
\item Решите неравенство. После умножения на знаменатель обязательно запишите выражения в скобках:
\[\frac{3(2-x)}5+\frac{10}5<\frac{7x}5-\frac{x+2}5.\]
\item Решите $3(2x+1)-5<16$. Найдите множество решений, принадлежащих отрезку $[-4;7]$, затем перечислите все целые решения из этого множества.
\item За перевод денег взимают $10$ руб. и $5\%$ от суммы перевода. На перевод вместе с комиссией есть $1000$ руб. Какую наибольшую целую сумму в рублях можно перевести?
\item Сумма трёх последовательных натуральных чётных чисел не больше $208$. Найдите наибольшую возможную тройку.
\item Найдите множество решений каждого неравенства:
\[
\text{а) }5x-3>5x+1;\quad
\text{б) }7x+2\le7x+3;\quad
\text{в) }2x+1<4x+5.
\]
\item Рассмотрите неравенство $3y+2a\le3y+8$, где $a$ --- параметр. При каких значениях $a$ подходят все действительные $y$? При каких значениях $a$ решений нет?
\end{enumerate}

\clearpage
\lessonsection{Ответы к тренировочному блоку}
\renewcommand{\arraystretch}{1.55}
\begin{tabularx}{\linewidth}{@{}p{1cm}X@{}}
\toprule
\textbf{№}&\textbf{Ответ}\\
\midrule
1 & $f(x)>0$ при $x<2$; $f(x)=0$ при $x=2$; $f(x)<0$ при $x>2$.\\
2 & $x\ge2$.\\
3 & $x\le-3$.\\
4 & $x\ge12$.\\
5 & $x>2$.\\
6 & $[-4;3)$; целые: $-4,-3,-2,-1,0,1,2$.\\
7 & $942$ руб.\\
8 & $66,68,70$.\\
9 & а) $\varnothing$; б) $x\in\mathbb R$; в) $x>-2$.\\
10 & При $a\le4$ --- все $y\in\mathbb R$; при $a>4$ решений нет.\\
\bottomrule
\end{tabularx}

\vspace{8mm}
\subsection*{Итоговая памятка}
\renewcommand{\arraystretch}{1.3}
\begin{tabularx}{\linewidth}{@{}p{.37\linewidth}X@{}}
\toprule
\textbf{Ситуация}&\textbf{Как действовать}\\
\midrule
Нужен знак функции & Смотрите, выше или ниже оси $Ox$ расположен график.\\
Сравниваем $f(x)$ с числом & Проведите горизонталь $y=c$ и считайте подходящие $x$.\\
Делим на отрицательное & Обязательно меняйте знак неравенства.\\
Вычитаем выражение & Заключите вычитаемое в скобки.\\
Дан дополнительный интервал & Найдите пересечение множеств решений.\\
Спрашивают целые числа & Учитывайте и границы, и целочисленность.\\
Переменная исчезла & Проверьте, истинно ли оставшееся числовое неравенство.\\
\bottomrule
\end{tabularx}
\vfill
\begin{center}
\small\color{muted}Лёвин Артём Александрович эксклюзивно для Екатерины Скелли
\end{center}
\end{document}
